% Separate research-scope record after CP32. No new full-group assertion.
\subsection*{The three-component marked boundary does not supply a large general-linear block}

Let \((A,\Theta)\) be a nonordinary principally polarized surface in
characteristic two, with two-rank \(f\in\{0,1\}\). Put
\(L=\mathcal O_A(3\Theta)\), and let \(C_i\) be three sharp smooth
odd-degree members in opposite total Picard twists, of degrees
\(a_1,a_2,a_3\), meeting transversely with no triple points. Write
\(d=a_1+a_2+a_3\), \(R=9d^2-1\),
\(r_i=9a_i^2-1\), and
\[
 \delta=18\sum_{i<j}a_i a_j.
\]
The graph of their connected union \(Y=C_1\cup C_2\cup C_3\)
has three vertices and \(\delta\) edges, so its cohomology has
dimension \(\delta-2\). The normalization sequence, after quotienting
the diagonal ambient cohomology, gives
\[
 0\longrightarrow T\longrightarrow S_Y\longrightarrow
 \frac{\bigoplus_{i=1}^3 H^1_{\mathrm{et}}(C_i,\mathbf F_2)}
      {H^1_{\mathrm{et}}(A,\mathbf F_2)}
 \longrightarrow0,
 \qquad \dim T=\delta-2.
\]
The ambient restriction is injective into each component. In particular
its image meets the graph kernel trivially, so the displayed sequence
is the actual relative sequence. Its dimension is
\[
 \dim S_Y=(\delta-2)+\sum_i(r_i+f)-f
          =9d^2-5+2f=R-4+2f.
\]
Thus the special unit space has the stated codimension four for
\(f=0\), or two for \(f=1\), in the generic odd-degree module.

On the family with these three normalization components marked,
\(T\) is a canonical sub-local-system. The component monodromies,
node permutations, and gluing translations supplied by the retained
boundary proofs all preserve it. Both \(T\) and its quotient are
nonzero. Consequently those family operations cannot act as
\(\mathrm{GL}(S_Y)\): the natural module of that full general-linear
group has no such proper invariant subspace. This is a statement
about the actual marked boundary representation, not about all
possible ambient monodromies obtained after leaving that family.

The canonical two-component addition proposition proves independent
natural blocks together with translations into an intersection heart.
Its later full-group theorem uses additional many-theta partition
comparisons in an even target where the relevant union already has
maximal unit rank. It does not identify a full general-linear action
on the larger special space \(S_Y\) in the odd target above. The
retained odd-plus-even obstruction note discusses even an actual
\(\mathrm{GL}_{R-6}\) block conditionally; it is not a proof of the
stronger \(\mathrm{GL}_{R-4+2f}\) premise.

In particular, the new full hyperplane-radical inertia theorem cannot
yet be combined with an asserted large boundary general-linear block
to bound its normal closure's fixed space by four or two dimensions.
That argument would first need actual additional boundary comparisons
and the positions of their submodules, with the odd target's rank drop
treated explicitly. No such comparison or fixed-space bound is claimed
in this record.

\paragraph{Read and verification scope.}
The entire canonical files
\path{agents/monodromy/nonordinary_ppav_odd_natural_blocks.tex},
\path{agents/monodromy/nonordinary_ppav_even_addition.tex}, and
\path{agents/monodromy/nonordinary_three_theta_graph_reduction.tex}
were read, together with the exact retained obstruction
\path{agents/monodromy/post_CP24_odd_even_boundary_obstruction.tex}
and \path{sections/A_route_limitations_29.tex}. The normalization
sequence and dimensions above were independently rederived. No new
numerical computation or primary-source reading is claimed. This
records an unavailable route premise; it does not refute full
monodromy for any geometric family.
